% Rev 7 -- Qualcomm Power & Limits Software Engineer
\documentclass[10pt,letterpaper]{article}

\usepackage[margin=0.50in]{geometry}
\usepackage[T1]{fontenc}
\usepackage[utf8]{inputenc}
\usepackage[default,type1]{sourcesanspro}
\usepackage[final]{microtype}
\usepackage[explicit]{titlesec}
\usepackage{enumitem}
\usepackage{tabularx}
\usepackage{array}
\usepackage{ragged2e}
\usepackage{iftex}
\usepackage[hidelinks]{hyperref}

\ifPDFTeX
  \input{glyphtounicode}
  \pdfgentounicode=1
\fi

\hypersetup{
  pdftitle={Connor Savugot Resume},
  pdfauthor={Connor Savugot},
  pdfsubject={Qualcomm Power and Limits Software Engineer},
  pdfkeywords={embedded software, embedded C, C++, Python, real-time systems, RTOS, FreeRTOS, bare metal, power control, power management, DSP, dsPIC33CK, PIC32, TMS320, JTAG, hardware software integration, multiprocessor systems, SoC, ARM, Embedded Linux, firmware validation, compiler toolchains, Git}
}

\pagestyle{empty}
\setcounter{secnumdepth}{0}
\setlength{\parindent}{0pt}
\setlength{\parskip}{0pt}
\hfuzz=0.4pt

\titleformat{\section}
  {\fontsize{11.3}{11.9}\selectfont\bfseries\uppercase}
  {}{0pt}{#1\hspace{0.12cm}\titlerule[0.65pt]}
\titlespacing{\section}{0pt}{0.14cm}{0.05cm}

\newlist{resumebullets}{itemize}{1}
\setlist[resumebullets]{
  label=\textbullet,
  leftmargin=0.42cm,
  topsep=0.004cm,
  itemsep=0.004cm,
  parsep=0pt,
  partopsep=0pt
}

\newcommand{\companyheader}[1]{%
  {\fontsize{10.4}{10.8}\selectfont\bfseries #1}\par\vspace{-0.03cm}
}

\newcommand{\resumeentry}[2]{%
  \begin{tabularx}{\textwidth}{@{}X>{\raggedleft\arraybackslash}p{2.40in}@{}}
    \textbf{#1} & #2
  \end{tabularx}\vspace{-0.075cm}
}

\newcommand{\roleentry}[2]{%
  \begin{tabularx}{\textwidth}{@{}X>{\raggedleft\arraybackslash}p{2.40in}@{}}
    \textbf{#1} & #2
  \end{tabularx}\vspace{-0.075cm}
}

\newcommand{\descriptorentry}[1]{%
  \begin{tabularx}{\textwidth}{@{}X>{\raggedleft\arraybackslash}p{2.40in}@{}}
    \textit{#1} &
  \end{tabularx}\vspace{-0.075cm}
}

\begin{document}
\fontsize{9.65}{10.50}\selectfont
\setlength{\RaggedRightRightskip}{0pt plus 5em}
\RaggedRight
\emergencystretch=1em
\hyphenpenalty=10000
\exhyphenpenalty=10000

\begin{center}
  {\fontsize{21.5}{22.2}\selectfont\bfseries Connor Savugot}\\[1.5pt]
  {\bfseries B.S. in Computer Engineering | North Carolina State University | Graduated May 2026}\\[1.5pt]
  {\fontsize{9.6}{10.2}\selectfont
    \href{mailto:consav04@gmail.com}{consav04@gmail.com}
    \enspace|\enspace \href{tel:+18285419487}{+1-828-541-9487}
    \enspace|\enspace Murphy, NC
    \enspace|\enspace \href{https://linkedin.com/in/connorsavugot}{linkedin.com/in/connorsavugot}
    \enspace|\enspace \href{https://github.com/Clem085}{github.com/Clem085}}
\end{center}
\vspace{-0.18cm}

\section{Technical Profile}
Embedded software engineer developing real-time C firmware for power, control, and communications systems across 16- and 32-bit microcontrollers, FreeRTOS, and bare-metal DSPs. Experience spans C++/Python engineering tools, multiprocessor hardware/software integration, datasheet-driven peripheral bring-up, JTAG/bench debugging, and Arm-based Embedded Linux systems.

\section{Technical Skills}
\textbf{Languages:} Embedded C, C++, Python, Bash\\
\textbf{Embedded \& Real-Time:} Bare-metal firmware, FreeRTOS, POSIX threads, interrupts/timers, peripheral/HAL bring-up, boot/update/recovery systems\\
\textbf{Processors:} dsPIC33CK (16-bit DSC), PIC32, TMS320 DSP, STMicroelectronics MCU, MSP430FR2355; TI AM62 Arm-based SoC\\
\textbf{Debug \& Build:} JTAG, register inspection, Code Composer Studio, GCC/Make, compiler/toolchain configuration, Git/GitLab, Docker, oscilloscopes, logic/protocol analyzers\\
\textbf{Systems \& Interfaces:} Embedded Linux, Yocto/Arago, device trees, systemd; PWM, ADC, I2C, SPI, UART/RS-232, CAN 2.0/CAN FD/J1939, IPMI/IPMB

\section{Engineering Experience}
\companyheader{AEGIS Power Systems}
\roleentry{Embedded Firmware Engineer}{May 2026--Present}
\begin{resumebullets}
  \item Develop low-level C firmware for 16-bit dsPIC33CK DSCs and bring up peripherals one subsystem at a time from datasheets, reference manuals, and schematics; isolate register, timing, configuration, and electrical faults using JTAG, targeted tests, oscilloscopes, and logic analyzers.
  \item Implemented and validated VITA 46.11 Tier 2 management firmware for a VPX power supply, exposing IPMI 2.0 sensor, FRU/inventory, and event data over an I2C-based IPMB link to the chassis manager.
  \item Integrate battery hardware and CAN-connected power devices with a TI AM62 Arm-based Embedded Linux controller; build Yocto/Arago images and kernels, update device trees, develop external-watchdog services, and trace faults across power, wiring, CAN, boot, configuration, and code.
\end{resumebullets}

\roleentry{Embedded Systems Intern}{May--Aug 2025}
\roleentry{Computer Science Intern}{May--Aug 2024}
\descriptorentry{Selected work across both internships}
\begin{resumebullets}
  \item Developed bare-metal C power-control firmware on a TMS320 DSP for transformer-coupled bidirectional buck--boost supplies using 12 PWM A/B pairs (24 physical outputs); configured each A channel's period/frequency, duty cycle, and relative phase while each B output followed or inverted its paired A signal, then validated ADC/PWM behavior in Code Composer Studio.
  \item Developed embedded C firmware across a two-processor power system: a PIC32 running FreeRTOS coordinated I2C, SPI, CAN/J1939, UART, and update behavior while the bare-metal TMS320 DSP executed real-time converter control.
  \item Built a resilient update/recovery workflow: a C++ CLI accepted HEX images, a Python TCP sender handled chunking and checksums, and the PIC32 staged new firmware and a golden image in external SPI flash, programmed and verified the DSP over UART, and controlled boot.
  \item In 2024, built a restricted SSH/chroot interface and Python/Bash commands for a resource-constrained TI Yocto/Arago power system; reverse-engineered RS-232 commands/checksums for programmable loads and relay controllers and helped assemble the associated test circuit.
\end{resumebullets}

\section{Selected Power \& Real-Time Projects}
\resumeentry{EV Active Sensor Adapter | Battery Monitoring \& CAN}{NC State University}
\begin{resumebullets}
  \item Helped architect a custom-PCB system integrating an STMicroelectronics MCU, TI BQ-family battery monitor/ADC, CAN transceivers, power/sensing interfaces, embedded control, and telemetry.
  \item Translated electrical and protocol constraints plus component trade studies into interface requirements, validation criteria, and staged subsystem tests; coordinated risks, design reviews, and documentation across the team.
\end{resumebullets}

\resumeentry{Real-Time Systems \& Formal Verification | C, Linux, Lustre}{Spring 2026}
\begin{resumebullets}
  \item Implemented and instrumented periodic POSIX-thread tasks in C using CPU affinity, absolute-time releases, priority-inheritance mutexes, and timestamp traces; analyzed EDF plus rate- and deadline-monotonic schedulability and plotted execution timelines.
  \item Modeled synchronous control logic in Lustre and used bounded model checking and counterexample traces to correct temporal monitors and verify required system invariants.
\end{resumebullets}

\resumeentry{Software Performance \& Machine Learning | C++, Python}{Selected Coursework}
\begin{resumebullets}
  \item Implemented BFS, DFS, and greedy best-first maze solvers in C++ and compared runtime and memory tradeoffs across maze configurations.
  \item Developed a Python mentor-matching pipeline using sentence-transformer embeddings and configurable weighted scoring; added CSV ingestion, persistent constraints/state, and explainable score output.
\end{resumebullets}

\section{Technical Leadership}
\resumeentry{ECE 306 Embedded Systems | Teaching Assistant \& Course Mentor}{Spring 2025--Spring 2026}
\begin{resumebullets}
  \item Delivered supplemental lectures and debugged embedded C/MSP430FR2355 firmware and custom-PCB projects across three semesters, diagnosing register, timing, peripheral, toolchain, wiring, and soldering failures across diverse student implementations.
\end{resumebullets}

\end{document}
